
\documentclass{article}
\usepackage{colortbl}
\usepackage{makecell}
\usepackage{multirow}
\usepackage{supertabular}

\begin{document}

\newcounter{utterance}

\centering \large Interaction Transcript for game `codenames', experiment `opponent\_easy', episode 2 with Qwen3.6{-}35B{-}A3B{-}FP8{-}without{-}reasoning.
\vspace{24pt}

{ \footnotesize  \setcounter{utterance}{1}
\setlength{\tabcolsep}{0pt}
\begin{supertabular}{c@{$\;$}|p{.15\linewidth}@{}p{.15\linewidth}p{.15\linewidth}p{.15\linewidth}p{.15\linewidth}p{.15\linewidth}}
    \# & $\;$A & \multicolumn{4}{c}{Game Master} & $\;\:$B\\
    \hline

    \theutterance \stepcounter{utterance}  
    & & & \multicolumn{2}{p{0.3\linewidth}}{
        \cellcolor[rgb]{0.9,0.9,0.9}{
            \makecell[{{p{\linewidth}}}]{
                \texttt{\tiny{[GM$|$GM]}}
                {'hidden': {'team': ['danza', 'prensa', 'enterrador', 'pez', 'cobre', 'washington', 'francia', 'regazo', 'corneta'], 'innocent': ['helicóptero', 'aceite', 'australia', 'cara', 'chocolate', 'agua', 'londres'], 'opponent': ['mosca', 'conductor', 'ornitorrinco', 'olympus', 'ballena', 'pasta', 'parque', 'ciclo'], 'assassin': ['enfermedad']}, 'revealed': {'team': {'team': [], 'innocent': [], 'opponent': [], 'assassin': []}, 'opponent': {'team': [], 'innocent': [], 'opponent': [], 'assassin': []}}}
            }
        }
    }
    & & \\ \\

    \theutterance \stepcounter{utterance}  
    & & \multicolumn{4}{p{0.6\linewidth}}{
        \cellcolor[rgb]{0.9,0.9,0.9}{
            \makecell[{{p{\linewidth}}}]{
                \texttt{\tiny{[P1$\langle$GM]}}
                \texttt{Proporciona una pista de una sola palabra que esté relacionada con algunas de las palabras de tu equipo de la lista a continuación. Puedes elegir con cuántas y con qué palabras se relaciona tu pista.} \\
\texttt{El objetivo es que tu compañero adivine todas las palabras de tu equipo basándose en las pistas que le des, antes de que el otro equipo haya descubierto todas las suyas. Si tu compañero adivina una palabra 'oponente' o 'distractor', su turno de adivinar termina y el equipo contrario puede hacer un turno. Si tu compañero adivina una palabra 'asesina', el juego termina y tu equipo pierde inmediatamente.} \\
\texttt{Cuantas más palabras relacione tu pista, más rápido podrás ganar el juego, pero más difíciles de adivinar serán las asociaciones para tu compañero, así que elige bien tu palabra pista. La palabra pista tiene que estar semánticamente relacionada con las palabras objetivo, no puede ser una de las palabras de las listas ni contener partes de ellas.} \\
\\ 
\texttt{Siempre da tu pista de una sola palabra y tu lista separada por comas de palabras objetivo relacionadas en el siguiente formato y haz tus respuestas lo más breves posible, nunca incluyas más texto que el requerido en este formato:} \\
\\ 
\texttt{PISTA: <PALABRA>} \\
\texttt{OBJETIVOS: <OBJETIVOS>} \\
\\ 
\texttt{Las palabras de tu equipo son: danza, prensa, enterrador, pez, cobre, washington, francia, regazo, corneta.} \\
\texttt{Las palabras del equipo de tu oponente son: mosca, conductor, ornitorrinco, olympus, ballena, pasta, parque, ciclo.} \\
\texttt{Las palabras distractor son: helicóptero, aceite, australia, cara, chocolate, agua, londres.} \\
\texttt{Las palabras asesinas son: enfermedad.} \\
            }
        }
    }
    & \\ \\

    \theutterance \stepcounter{utterance}  
    & \multicolumn{4}{p{0.6\linewidth}}{
        \cellcolor[rgb]{0.9,0.9,0.9}{
            \makecell[{{p{\linewidth}}}]{
                \texttt{\tiny{[P1$\rangle$GM]}}
                \texttt{PISTA: música} \\
\texttt{OBJETIVOS: danza, corneta} \\
            }
        }
    }
    & & \\ \\

    \theutterance \stepcounter{utterance}  
    & & & \multicolumn{2}{p{0.3\linewidth}}{
        \cellcolor[rgb]{0.9,0.9,0.9}{
            \makecell[{{p{\linewidth}}}]{
                \texttt{\tiny{[GM$|$GM]}}
                \texttt{música} \\
            }
        }
    }
    & & \\ \\

    \theutterance \stepcounter{utterance}  
    & & & \multicolumn{2}{p{0.3\linewidth}}{
        \cellcolor[rgb]{0.9,0.9,0.9}{
            \makecell[{{p{\linewidth}}}]{
                \texttt{\tiny{[GM$|$GM]}}
                ['danza', 'corneta']
            }
        }
    }
    & & \\ \\

    \theutterance \stepcounter{utterance}  
    & & & \multicolumn{2}{p{0.3\linewidth}}{
        \cellcolor[rgb]{0.9,0.9,0.9}{
            \makecell[{{p{\linewidth}}}]{
                \texttt{\tiny{[GM$|$GM]}}
                {'word': 'danza', 'assignment': 'team'}
            }
        }
    }
    & & \\ \\

    \theutterance \stepcounter{utterance}  
    & & & \multicolumn{2}{p{0.3\linewidth}}{
        \cellcolor[rgb]{0.9,0.9,0.9}{
            \makecell[{{p{\linewidth}}}]{
                \texttt{\tiny{[GM$|$GM]}}
                {'word': 'corneta', 'assignment': 'team'}
            }
        }
    }
    & & \\ \\

    \theutterance \stepcounter{utterance}  
    & & \multicolumn{4}{p{0.6\linewidth}}{
        \cellcolor[rgb]{0.9,0.9,0.9}{
            \makecell[{{p{\linewidth}}}]{
                \texttt{\tiny{[GM$\rangle$P2]}}
                \texttt{Proporciona una lista separada por comas de hasta 2 palabras de la siguiente lista que mejor se relacionen o estén más estrechamente asociadas con la palabra 'música'. Comienza siempre tu lista de suposición(es) con ADIVINANZA:  y no incluyas ningún otro texto en tu respuesta.} \\
\\ 
\texttt{olympus, mosca, washington, australia, enterrador, corneta, danza, aceite, pez, cobre, regazo, enfermedad, pasta, ciclo, helicóptero, agua, chocolate, prensa, parque, francia, conductor, ballena, ornitorrinco, londres, cara} \\
            }
        }
    }
    & \\ \\

    \theutterance \stepcounter{utterance}  
    & & & \multicolumn{4}{p{0.6\linewidth}}{
        \cellcolor[rgb]{0.9,0.9,0.9}{
            \makecell[{{p{\linewidth}}}]{
                \texttt{\tiny{[GM$\langle$P2]}}
                \texttt{ADIVINANZA: corneta, danza} \\
            }
        }
    }
     \\ \\

    \theutterance \stepcounter{utterance}  
    & & & \multicolumn{2}{p{0.3\linewidth}}{
        \cellcolor[rgb]{0.9,0.9,0.9}{
            \makecell[{{p{\linewidth}}}]{
                \texttt{\tiny{[GM$|$GM]}}
                ['corneta', 'danza']
            }
        }
    }
    & & \\ \\

    \theutterance \stepcounter{utterance}  
    & & & \multicolumn{2}{p{0.3\linewidth}}{
        \cellcolor[rgb]{0.9,0.9,0.9}{
            \makecell[{{p{\linewidth}}}]{
                \texttt{\tiny{[GM$|$GM]}}
                {'word': 'corneta', 'assignment': 'team'}
            }
        }
    }
    & & \\ \\

    \theutterance \stepcounter{utterance}  
    & & & \multicolumn{2}{p{0.3\linewidth}}{
        \cellcolor[rgb]{0.9,0.9,0.9}{
            \makecell[{{p{\linewidth}}}]{
                \texttt{\tiny{[GM$|$GM]}}
                {'word': 'corneta', 'assignment': 'team'}
            }
        }
    }
    & & \\ \\

    \theutterance \stepcounter{utterance}  
    & & & \multicolumn{2}{p{0.3\linewidth}}{
        \cellcolor[rgb]{0.9,0.9,0.9}{
            \makecell[{{p{\linewidth}}}]{
                \texttt{\tiny{[GM$|$GM]}}
                {'word': 'danza', 'assignment': 'team'}
            }
        }
    }
    & & \\ \\

    \theutterance \stepcounter{utterance}  
    & & & \multicolumn{2}{p{0.3\linewidth}}{
        \cellcolor[rgb]{0.9,0.9,0.9}{
            \makecell[{{p{\linewidth}}}]{
                \texttt{\tiny{[GM$|$GM]}}
                {'word': 'danza', 'assignment': 'team'}
            }
        }
    }
    & & \\ \\

    \theutterance \stepcounter{utterance}  
    & & & \multicolumn{2}{p{0.3\linewidth}}{
        \cellcolor[rgb]{0.9,0.9,0.9}{
            \makecell[{{p{\linewidth}}}]{
                \texttt{\tiny{[GM$|$GM]}}
                {'word': 'conductor', 'assignment': 'opponent'}
            }
        }
    }
    & & \\ \\

    \theutterance \stepcounter{utterance}  
    & & & \multicolumn{2}{p{0.3\linewidth}}{
        \cellcolor[rgb]{0.9,0.9,0.9}{
            \makecell[{{p{\linewidth}}}]{
                \texttt{\tiny{[GM$|$GM]}}
                {'hidden': {'team': ['prensa', 'enterrador', 'pez', 'cobre', 'washington', 'francia', 'regazo'], 'innocent': ['helicóptero', 'aceite', 'australia', 'cara', 'chocolate', 'agua', 'londres'], 'opponent': ['mosca', 'ornitorrinco', 'olympus', 'ballena', 'pasta', 'parque', 'ciclo'], 'assassin': ['enfermedad']}, 'revealed': {'team': {'team': ['corneta', 'danza'], 'innocent': [], 'opponent': [], 'assassin': []}, 'opponent': {'team': [], 'innocent': [], 'opponent': ['conductor'], 'assassin': []}}}
            }
        }
    }
    & & \\ \\

    \theutterance \stepcounter{utterance}  
    & & \multicolumn{4}{p{0.6\linewidth}}{
        \cellcolor[rgb]{0.9,0.9,0.9}{
            \makecell[{{p{\linewidth}}}]{
                \texttt{\tiny{[P1$\langle$GM]}}
                \texttt{Las palabras corneta, danza fueron adivinadas correctamente. El turno de tu compañero terminó allí.} \\
\texttt{Ahora proporcione otra pista relacionada con algunas de las palabras restantes de su equipo y una lista de las palabras objetivo relacionadas. Recuerde que debe comenzar su pista con 'PISTA: ', poner una nueva línea, y comenzar su lista de palabras objetivo separadas por comas con 'OBJETIVOS: '. Aviso: se han eliminado algunas palabras de las listas con respecto a solicitudes anteriores.} \\
\\ 
\texttt{Las palabras restantes de su equipo son: prensa, enterrador, pez, cobre, washington, francia, regazo.} \\
\texttt{Las palabras restantes para su adversario son: mosca, conductor, ornitorrinco, olympus, ballena, pasta, parque, ciclo.} \\
\texttt{Las palabras distractoras restantes son: helicóptero, aceite, australia, cara, chocolate, agua, londres.} \\
\texttt{Las palabras asesinas restantes son: enfermedad.} \\
            }
        }
    }
    & \\ \\

    \theutterance \stepcounter{utterance}  
    & \multicolumn{4}{p{0.6\linewidth}}{
        \cellcolor[rgb]{0.9,0.9,0.9}{
            \makecell[{{p{\linewidth}}}]{
                \texttt{\tiny{[P1$\rangle$GM]}}
                \texttt{PISTA: metal} \\
\texttt{OBJETIVOS: cobre} \\
            }
        }
    }
    & & \\ \\

    \theutterance \stepcounter{utterance}  
    & & & \multicolumn{2}{p{0.3\linewidth}}{
        \cellcolor[rgb]{0.9,0.9,0.9}{
            \makecell[{{p{\linewidth}}}]{
                \texttt{\tiny{[GM$|$GM]}}
                \texttt{metal} \\
            }
        }
    }
    & & \\ \\

    \theutterance \stepcounter{utterance}  
    & & & \multicolumn{2}{p{0.3\linewidth}}{
        \cellcolor[rgb]{0.9,0.9,0.9}{
            \makecell[{{p{\linewidth}}}]{
                \texttt{\tiny{[GM$|$GM]}}
                ['cobre']
            }
        }
    }
    & & \\ \\

    \theutterance \stepcounter{utterance}  
    & & & \multicolumn{2}{p{0.3\linewidth}}{
        \cellcolor[rgb]{0.9,0.9,0.9}{
            \makecell[{{p{\linewidth}}}]{
                \texttt{\tiny{[GM$|$GM]}}
                {'word': 'cobre', 'assignment': 'team'}
            }
        }
    }
    & & \\ \\

    \theutterance \stepcounter{utterance}  
    & & \multicolumn{4}{p{0.6\linewidth}}{
        \cellcolor[rgb]{0.9,0.9,0.9}{
            \makecell[{{p{\linewidth}}}]{
                \texttt{\tiny{[GM$\rangle$P2]}}
                \texttt{Las palabras corneta, danza fueron adivinadas correctamente. Tu turno terminó allí.} \\
\texttt{Ahora proporcione otra lista separada por comas de al menos 1 y hasta 1 palabras de la siguiente lista de palabras que mejor se relacionen o estén más estrechamente asociadas con la palabra 'metal'. Recuerde comenzar su respuesta con ADIVINANZA: . Tenga en cuenta: algunas palabras se han eliminado de la lista en comparación con solicitudes anteriores.} \\
\\ 
\texttt{olympus, mosca, washington, australia, enterrador, aceite, pez, cobre, regazo, enfermedad, pasta, ciclo, helicóptero, agua, chocolate, prensa, parque, francia, ballena, ornitorrinco, londres, cara} \\
            }
        }
    }
    & \\ \\

    \theutterance \stepcounter{utterance}  
    & & & \multicolumn{4}{p{0.6\linewidth}}{
        \cellcolor[rgb]{0.9,0.9,0.9}{
            \makecell[{{p{\linewidth}}}]{
                \texttt{\tiny{[GM$\langle$P2]}}
                \texttt{ADIVINANZA: cobre} \\
            }
        }
    }
     \\ \\

    \theutterance \stepcounter{utterance}  
    & & & \multicolumn{2}{p{0.3\linewidth}}{
        \cellcolor[rgb]{0.9,0.9,0.9}{
            \makecell[{{p{\linewidth}}}]{
                \texttt{\tiny{[GM$|$GM]}}
                ['cobre']
            }
        }
    }
    & & \\ \\

    \theutterance \stepcounter{utterance}  
    & & & \multicolumn{2}{p{0.3\linewidth}}{
        \cellcolor[rgb]{0.9,0.9,0.9}{
            \makecell[{{p{\linewidth}}}]{
                \texttt{\tiny{[GM$|$GM]}}
                {'word': 'cobre', 'assignment': 'team'}
            }
        }
    }
    & & \\ \\

    \theutterance \stepcounter{utterance}  
    & & & \multicolumn{2}{p{0.3\linewidth}}{
        \cellcolor[rgb]{0.9,0.9,0.9}{
            \makecell[{{p{\linewidth}}}]{
                \texttt{\tiny{[GM$|$GM]}}
                {'word': 'cobre', 'assignment': 'team'}
            }
        }
    }
    & & \\ \\

    \theutterance \stepcounter{utterance}  
    & & & \multicolumn{2}{p{0.3\linewidth}}{
        \cellcolor[rgb]{0.9,0.9,0.9}{
            \makecell[{{p{\linewidth}}}]{
                \texttt{\tiny{[GM$|$GM]}}
                {'word': 'ornitorrinco', 'assignment': 'opponent'}
            }
        }
    }
    & & \\ \\

    \theutterance \stepcounter{utterance}  
    & & & \multicolumn{2}{p{0.3\linewidth}}{
        \cellcolor[rgb]{0.9,0.9,0.9}{
            \makecell[{{p{\linewidth}}}]{
                \texttt{\tiny{[GM$|$GM]}}
                {'hidden': {'team': ['prensa', 'enterrador', 'pez', 'washington', 'francia', 'regazo'], 'innocent': ['helicóptero', 'aceite', 'australia', 'cara', 'chocolate', 'agua', 'londres'], 'opponent': ['mosca', 'olympus', 'ballena', 'pasta', 'parque', 'ciclo'], 'assassin': ['enfermedad']}, 'revealed': {'team': {'team': ['corneta', 'danza', 'cobre'], 'innocent': [], 'opponent': [], 'assassin': []}, 'opponent': {'team': [], 'innocent': [], 'opponent': ['conductor', 'ornitorrinco'], 'assassin': []}}}
            }
        }
    }
    & & \\ \\

    \theutterance \stepcounter{utterance}  
    & & \multicolumn{4}{p{0.6\linewidth}}{
        \cellcolor[rgb]{0.9,0.9,0.9}{
            \makecell[{{p{\linewidth}}}]{
                \texttt{\tiny{[P1$\langle$GM]}}
                \texttt{La palabra cobre fue adivinada correctamente. El turno de tu compañero terminó allí.} \\
\texttt{Ahora proporcione otra pista relacionada con algunas de las palabras restantes de su equipo y una lista de las palabras objetivo relacionadas. Recuerde que debe comenzar su pista con 'PISTA: ', poner una nueva línea, y comenzar su lista de palabras objetivo separadas por comas con 'OBJETIVOS: '. Aviso: se han eliminado algunas palabras de las listas con respecto a solicitudes anteriores.} \\
\\ 
\texttt{Las palabras restantes de su equipo son: prensa, enterrador, pez, washington, francia, regazo.} \\
\texttt{Las palabras restantes para su adversario son: mosca, ornitorrinco, olympus, ballena, pasta, parque, ciclo.} \\
\texttt{Las palabras distractoras restantes son: helicóptero, aceite, australia, cara, chocolate, agua, londres.} \\
\texttt{Las palabras asesinas restantes son: enfermedad.} \\
            }
        }
    }
    & \\ \\

    \theutterance \stepcounter{utterance}  
    & \multicolumn{4}{p{0.6\linewidth}}{
        \cellcolor[rgb]{0.9,0.9,0.9}{
            \makecell[{{p{\linewidth}}}]{
                \texttt{\tiny{[P1$\rangle$GM]}}
                \texttt{PISTA: país} \\
\texttt{OBJETIVOS: francia} \\
            }
        }
    }
    & & \\ \\

    \theutterance \stepcounter{utterance}  
    & & & \multicolumn{2}{p{0.3\linewidth}}{
        \cellcolor[rgb]{0.9,0.9,0.9}{
            \makecell[{{p{\linewidth}}}]{
                \texttt{\tiny{[GM$|$GM]}}
                \texttt{país} \\
            }
        }
    }
    & & \\ \\

    \theutterance \stepcounter{utterance}  
    & & & \multicolumn{2}{p{0.3\linewidth}}{
        \cellcolor[rgb]{0.9,0.9,0.9}{
            \makecell[{{p{\linewidth}}}]{
                \texttt{\tiny{[GM$|$GM]}}
                ['francia']
            }
        }
    }
    & & \\ \\

    \theutterance \stepcounter{utterance}  
    & & & \multicolumn{2}{p{0.3\linewidth}}{
        \cellcolor[rgb]{0.9,0.9,0.9}{
            \makecell[{{p{\linewidth}}}]{
                \texttt{\tiny{[GM$|$GM]}}
                {'word': 'francia', 'assignment': 'team'}
            }
        }
    }
    & & \\ \\

    \theutterance \stepcounter{utterance}  
    & & \multicolumn{4}{p{0.6\linewidth}}{
        \cellcolor[rgb]{0.9,0.9,0.9}{
            \makecell[{{p{\linewidth}}}]{
                \texttt{\tiny{[GM$\rangle$P2]}}
                \texttt{La palabra cobre fue adivinada correctamente. Tu turno terminó allí.} \\
\texttt{Ahora proporcione otra lista separada por comas de al menos 1 y hasta 1 palabras de la siguiente lista de palabras que mejor se relacionen o estén más estrechamente asociadas con la palabra 'país'. Recuerde comenzar su respuesta con ADIVINANZA: . Tenga en cuenta: algunas palabras se han eliminado de la lista en comparación con solicitudes anteriores.} \\
\\ 
\texttt{olympus, mosca, washington, australia, enterrador, aceite, pez, regazo, enfermedad, pasta, ciclo, helicóptero, agua, chocolate, prensa, parque, francia, ballena, londres, cara} \\
            }
        }
    }
    & \\ \\

    \theutterance \stepcounter{utterance}  
    & & & \multicolumn{4}{p{0.6\linewidth}}{
        \cellcolor[rgb]{0.9,0.9,0.9}{
            \makecell[{{p{\linewidth}}}]{
                \texttt{\tiny{[GM$\langle$P2]}}
                \texttt{ADIVINANZA: australia} \\
            }
        }
    }
     \\ \\

    \theutterance \stepcounter{utterance}  
    & & & \multicolumn{2}{p{0.3\linewidth}}{
        \cellcolor[rgb]{0.9,0.9,0.9}{
            \makecell[{{p{\linewidth}}}]{
                \texttt{\tiny{[GM$|$GM]}}
                ['australia']
            }
        }
    }
    & & \\ \\

    \theutterance \stepcounter{utterance}  
    & & & \multicolumn{2}{p{0.3\linewidth}}{
        \cellcolor[rgb]{0.9,0.9,0.9}{
            \makecell[{{p{\linewidth}}}]{
                \texttt{\tiny{[GM$|$GM]}}
                {'word': 'australia', 'assignment': 'innocent'}
            }
        }
    }
    & & \\ \\

    \theutterance \stepcounter{utterance}  
    & & & \multicolumn{2}{p{0.3\linewidth}}{
        \cellcolor[rgb]{0.9,0.9,0.9}{
            \makecell[{{p{\linewidth}}}]{
                \texttt{\tiny{[GM$|$GM]}}
                \texttt{australia} \\
            }
        }
    }
    & & \\ \\

    \theutterance \stepcounter{utterance}  
    & & & \multicolumn{2}{p{0.3\linewidth}}{
        \cellcolor[rgb]{0.9,0.9,0.9}{
            \makecell[{{p{\linewidth}}}]{
                \texttt{\tiny{[GM$|$GM]}}
                {'word': 'parque', 'assignment': 'opponent'}
            }
        }
    }
    & & \\ \\

    \theutterance \stepcounter{utterance}  
    & & & \multicolumn{2}{p{0.3\linewidth}}{
        \cellcolor[rgb]{0.9,0.9,0.9}{
            \makecell[{{p{\linewidth}}}]{
                \texttt{\tiny{[GM$|$GM]}}
                {'hidden': {'team': ['prensa', 'enterrador', 'pez', 'washington', 'francia', 'regazo'], 'innocent': ['helicóptero', 'aceite', 'cara', 'chocolate', 'agua', 'londres'], 'opponent': ['mosca', 'olympus', 'ballena', 'pasta', 'ciclo'], 'assassin': ['enfermedad']}, 'revealed': {'team': {'team': ['corneta', 'danza', 'cobre'], 'innocent': ['australia'], 'opponent': [], 'assassin': []}, 'opponent': {'team': [], 'innocent': [], 'opponent': ['conductor', 'ornitorrinco', 'parque'], 'assassin': []}}}
            }
        }
    }
    & & \\ \\

    \theutterance \stepcounter{utterance}  
    & & \multicolumn{4}{p{0.6\linewidth}}{
        \cellcolor[rgb]{0.9,0.9,0.9}{
            \makecell[{{p{\linewidth}}}]{
                \texttt{\tiny{[P1$\langle$GM]}}
                \texttt{La palabra australia fue adivinada, pero es una palabra inocente. El turno de tu compañero terminó allí.} \\
\texttt{Ahora proporcione otra pista relacionada con algunas de las palabras restantes de su equipo y una lista de las palabras objetivo relacionadas. Recuerde que debe comenzar su pista con 'PISTA: ', poner una nueva línea, y comenzar su lista de palabras objetivo separadas por comas con 'OBJETIVOS: '. Aviso: se han eliminado algunas palabras de las listas con respecto a solicitudes anteriores.} \\
\\ 
\texttt{Las palabras restantes de su equipo son: prensa, enterrador, pez, washington, francia, regazo.} \\
\texttt{Las palabras restantes para su adversario son: mosca, olympus, ballena, pasta, parque, ciclo.} \\
\texttt{Las palabras distractoras restantes son: helicóptero, aceite, cara, chocolate, agua, londres.} \\
\texttt{Las palabras asesinas restantes son: enfermedad.} \\
            }
        }
    }
    & \\ \\

    \theutterance \stepcounter{utterance}  
    & \multicolumn{4}{p{0.6\linewidth}}{
        \cellcolor[rgb]{0.9,0.9,0.9}{
            \makecell[{{p{\linewidth}}}]{
                \texttt{\tiny{[P1$\rangle$GM]}}
                \texttt{PISTA: tierra} \\
\texttt{OBJETIVOS: enterrador, pez} \\
            }
        }
    }
    & & \\ \\

    \theutterance \stepcounter{utterance}  
    & & & \multicolumn{2}{p{0.3\linewidth}}{
        \cellcolor[rgb]{0.9,0.9,0.9}{
            \makecell[{{p{\linewidth}}}]{
                \texttt{\tiny{[GM$|$GM]}}
                \texttt{tierra} \\
            }
        }
    }
    & & \\ \\

    \theutterance \stepcounter{utterance}  
    & & & \multicolumn{2}{p{0.3\linewidth}}{
        \cellcolor[rgb]{0.9,0.9,0.9}{
            \makecell[{{p{\linewidth}}}]{
                \texttt{\tiny{[GM$|$GM]}}
                ['enterrador', 'pez']
            }
        }
    }
    & & \\ \\

    \theutterance \stepcounter{utterance}  
    & & & \multicolumn{2}{p{0.3\linewidth}}{
        \cellcolor[rgb]{0.9,0.9,0.9}{
            \makecell[{{p{\linewidth}}}]{
                \texttt{\tiny{[GM$|$GM]}}
                {'word': 'enterrador', 'assignment': 'team'}
            }
        }
    }
    & & \\ \\

    \theutterance \stepcounter{utterance}  
    & & & \multicolumn{2}{p{0.3\linewidth}}{
        \cellcolor[rgb]{0.9,0.9,0.9}{
            \makecell[{{p{\linewidth}}}]{
                \texttt{\tiny{[GM$|$GM]}}
                {'word': 'pez', 'assignment': 'team'}
            }
        }
    }
    & & \\ \\

    \theutterance \stepcounter{utterance}  
    & & \multicolumn{4}{p{0.6\linewidth}}{
        \cellcolor[rgb]{0.9,0.9,0.9}{
            \makecell[{{p{\linewidth}}}]{
                \texttt{\tiny{[GM$\rangle$P2]}}
                \texttt{La palabra australia fue adivinada, pero es una palabra inocente. Tu turno terminó allí.} \\
\texttt{Ahora proporcione otra lista separada por comas de al menos 1 y hasta 2 palabras de la siguiente lista de palabras que mejor se relacionen o estén más estrechamente asociadas con la palabra 'tierra'. Recuerde comenzar su respuesta con ADIVINANZA: . Tenga en cuenta: algunas palabras se han eliminado de la lista en comparación con solicitudes anteriores.} \\
\\ 
\texttt{olympus, mosca, washington, enterrador, aceite, pez, regazo, enfermedad, pasta, ciclo, helicóptero, agua, chocolate, prensa, francia, ballena, londres, cara} \\
            }
        }
    }
    & \\ \\

    \theutterance \stepcounter{utterance}  
    & & & \multicolumn{4}{p{0.6\linewidth}}{
        \cellcolor[rgb]{0.9,0.9,0.9}{
            \makecell[{{p{\linewidth}}}]{
                \texttt{\tiny{[GM$\langle$P2]}}
                \texttt{ADIVINANZA: olympus, londres} \\
            }
        }
    }
     \\ \\

    \theutterance \stepcounter{utterance}  
    & & & \multicolumn{2}{p{0.3\linewidth}}{
        \cellcolor[rgb]{0.9,0.9,0.9}{
            \makecell[{{p{\linewidth}}}]{
                \texttt{\tiny{[GM$|$GM]}}
                ['olympus', 'londres']
            }
        }
    }
    & & \\ \\

    \theutterance \stepcounter{utterance}  
    & & & \multicolumn{2}{p{0.3\linewidth}}{
        \cellcolor[rgb]{0.9,0.9,0.9}{
            \makecell[{{p{\linewidth}}}]{
                \texttt{\tiny{[GM$|$GM]}}
                {'word': 'olympus', 'assignment': 'opponent'}
            }
        }
    }
    & & \\ \\

    \theutterance \stepcounter{utterance}  
    & & & \multicolumn{2}{p{0.3\linewidth}}{
        \cellcolor[rgb]{0.9,0.9,0.9}{
            \makecell[{{p{\linewidth}}}]{
                \texttt{\tiny{[GM$|$GM]}}
                \texttt{olympus} \\
            }
        }
    }
    & & \\ \\

    \theutterance \stepcounter{utterance}  
    & & & \multicolumn{2}{p{0.3\linewidth}}{
        \cellcolor[rgb]{0.9,0.9,0.9}{
            \makecell[{{p{\linewidth}}}]{
                \texttt{\tiny{[GM$|$GM]}}
                {'word': 'pasta', 'assignment': 'opponent'}
            }
        }
    }
    & & \\ \\

    \theutterance \stepcounter{utterance}  
    & & & \multicolumn{2}{p{0.3\linewidth}}{
        \cellcolor[rgb]{0.9,0.9,0.9}{
            \makecell[{{p{\linewidth}}}]{
                \texttt{\tiny{[GM$|$GM]}}
                {'hidden': {'team': ['prensa', 'enterrador', 'pez', 'washington', 'francia', 'regazo'], 'innocent': ['helicóptero', 'aceite', 'cara', 'chocolate', 'agua', 'londres'], 'opponent': ['mosca', 'ballena', 'ciclo'], 'assassin': ['enfermedad']}, 'revealed': {'team': {'team': ['corneta', 'danza', 'cobre'], 'innocent': ['australia'], 'opponent': ['olympus'], 'assassin': []}, 'opponent': {'team': [], 'innocent': [], 'opponent': ['conductor', 'ornitorrinco', 'parque', 'pasta'], 'assassin': []}}}
            }
        }
    }
    & & \\ \\

    \theutterance \stepcounter{utterance}  
    & & \multicolumn{4}{p{0.6\linewidth}}{
        \cellcolor[rgb]{0.9,0.9,0.9}{
            \makecell[{{p{\linewidth}}}]{
                \texttt{\tiny{[P1$\langle$GM]}}
                \texttt{La palabra olympus fue adivinada, pero es una palabra oponente. El turno de tu compañero terminó allí.} \\
\texttt{Ahora proporcione otra pista relacionada con algunas de las palabras restantes de su equipo y una lista de las palabras objetivo relacionadas. Recuerde que debe comenzar su pista con 'PISTA: ', poner una nueva línea, y comenzar su lista de palabras objetivo separadas por comas con 'OBJETIVOS: '. Aviso: se han eliminado algunas palabras de las listas con respecto a solicitudes anteriores.} \\
\\ 
\texttt{Las palabras restantes de su equipo son: prensa, enterrador, pez, washington, francia, regazo.} \\
\texttt{Las palabras restantes para su adversario son: mosca, ballena, pasta, ciclo.} \\
\texttt{Las palabras distractoras restantes son: helicóptero, aceite, cara, chocolate, agua, londres.} \\
\texttt{Las palabras asesinas restantes son: enfermedad.} \\
            }
        }
    }
    & \\ \\

    \theutterance \stepcounter{utterance}  
    & \multicolumn{4}{p{0.6\linewidth}}{
        \cellcolor[rgb]{0.9,0.9,0.9}{
            \makecell[{{p{\linewidth}}}]{
                \texttt{\tiny{[P1$\rangle$GM]}}
                \texttt{PISTA: estado} \\
\texttt{OBJETIVOS: washington} \\
            }
        }
    }
    & & \\ \\

    \theutterance \stepcounter{utterance}  
    & & & \multicolumn{2}{p{0.3\linewidth}}{
        \cellcolor[rgb]{0.9,0.9,0.9}{
            \makecell[{{p{\linewidth}}}]{
                \texttt{\tiny{[GM$|$GM]}}
                \texttt{estado} \\
            }
        }
    }
    & & \\ \\

    \theutterance \stepcounter{utterance}  
    & & & \multicolumn{2}{p{0.3\linewidth}}{
        \cellcolor[rgb]{0.9,0.9,0.9}{
            \makecell[{{p{\linewidth}}}]{
                \texttt{\tiny{[GM$|$GM]}}
                ['washington']
            }
        }
    }
    & & \\ \\

    \theutterance \stepcounter{utterance}  
    & & & \multicolumn{2}{p{0.3\linewidth}}{
        \cellcolor[rgb]{0.9,0.9,0.9}{
            \makecell[{{p{\linewidth}}}]{
                \texttt{\tiny{[GM$|$GM]}}
                {'word': 'washington', 'assignment': 'team'}
            }
        }
    }
    & & \\ \\

    \theutterance \stepcounter{utterance}  
    & & \multicolumn{4}{p{0.6\linewidth}}{
        \cellcolor[rgb]{0.9,0.9,0.9}{
            \makecell[{{p{\linewidth}}}]{
                \texttt{\tiny{[GM$\rangle$P2]}}
                \texttt{La palabra olympus fue adivinada, pero es una palabra oponente. Tu turno terminó allí.} \\
\texttt{Ahora proporcione otra lista separada por comas de al menos 1 y hasta 1 palabras de la siguiente lista de palabras que mejor se relacionen o estén más estrechamente asociadas con la palabra 'estado'. Recuerde comenzar su respuesta con ADIVINANZA: . Tenga en cuenta: algunas palabras se han eliminado de la lista en comparación con solicitudes anteriores.} \\
\\ 
\texttt{mosca, washington, enterrador, aceite, pez, regazo, enfermedad, ciclo, helicóptero, agua, chocolate, prensa, francia, ballena, londres, cara} \\
            }
        }
    }
    & \\ \\

    \theutterance \stepcounter{utterance}  
    & & & \multicolumn{4}{p{0.6\linewidth}}{
        \cellcolor[rgb]{0.9,0.9,0.9}{
            \makecell[{{p{\linewidth}}}]{
                \texttt{\tiny{[GM$\langle$P2]}}
                \texttt{ADIVINANZA: washington} \\
            }
        }
    }
     \\ \\

    \theutterance \stepcounter{utterance}  
    & & & \multicolumn{2}{p{0.3\linewidth}}{
        \cellcolor[rgb]{0.9,0.9,0.9}{
            \makecell[{{p{\linewidth}}}]{
                \texttt{\tiny{[GM$|$GM]}}
                ['washington']
            }
        }
    }
    & & \\ \\

    \theutterance \stepcounter{utterance}  
    & & & \multicolumn{2}{p{0.3\linewidth}}{
        \cellcolor[rgb]{0.9,0.9,0.9}{
            \makecell[{{p{\linewidth}}}]{
                \texttt{\tiny{[GM$|$GM]}}
                {'word': 'washington', 'assignment': 'team'}
            }
        }
    }
    & & \\ \\

    \theutterance \stepcounter{utterance}  
    & & & \multicolumn{2}{p{0.3\linewidth}}{
        \cellcolor[rgb]{0.9,0.9,0.9}{
            \makecell[{{p{\linewidth}}}]{
                \texttt{\tiny{[GM$|$GM]}}
                {'word': 'washington', 'assignment': 'team'}
            }
        }
    }
    & & \\ \\

    \theutterance \stepcounter{utterance}  
    & & & \multicolumn{2}{p{0.3\linewidth}}{
        \cellcolor[rgb]{0.9,0.9,0.9}{
            \makecell[{{p{\linewidth}}}]{
                \texttt{\tiny{[GM$|$GM]}}
                {'word': 'mosca', 'assignment': 'opponent'}
            }
        }
    }
    & & \\ \\

    \theutterance \stepcounter{utterance}  
    & & & \multicolumn{2}{p{0.3\linewidth}}{
        \cellcolor[rgb]{0.9,0.9,0.9}{
            \makecell[{{p{\linewidth}}}]{
                \texttt{\tiny{[GM$|$GM]}}
                {'hidden': {'team': ['prensa', 'enterrador', 'pez', 'francia', 'regazo'], 'innocent': ['helicóptero', 'aceite', 'cara', 'chocolate', 'agua', 'londres'], 'opponent': ['ballena', 'ciclo'], 'assassin': ['enfermedad']}, 'revealed': {'team': {'team': ['corneta', 'danza', 'cobre', 'washington'], 'innocent': ['australia'], 'opponent': ['olympus'], 'assassin': []}, 'opponent': {'team': [], 'innocent': [], 'opponent': ['conductor', 'ornitorrinco', 'parque', 'pasta', 'mosca'], 'assassin': []}}}
            }
        }
    }
    & & \\ \\

    \theutterance \stepcounter{utterance}  
    & & \multicolumn{4}{p{0.6\linewidth}}{
        \cellcolor[rgb]{0.9,0.9,0.9}{
            \makecell[{{p{\linewidth}}}]{
                \texttt{\tiny{[P1$\langle$GM]}}
                \texttt{La palabra washington fue adivinada correctamente. El turno de tu compañero terminó allí.} \\
\texttt{Ahora proporcione otra pista relacionada con algunas de las palabras restantes de su equipo y una lista de las palabras objetivo relacionadas. Recuerde que debe comenzar su pista con 'PISTA: ', poner una nueva línea, y comenzar su lista de palabras objetivo separadas por comas con 'OBJETIVOS: '. Aviso: se han eliminado algunas palabras de las listas con respecto a solicitudes anteriores.} \\
\\ 
\texttt{Las palabras restantes de su equipo son: prensa, enterrador, pez, francia, regazo.} \\
\texttt{Las palabras restantes para su adversario son: mosca, ballena, ciclo.} \\
\texttt{Las palabras distractoras restantes son: helicóptero, aceite, cara, chocolate, agua, londres.} \\
\texttt{Las palabras asesinas restantes son: enfermedad.} \\
            }
        }
    }
    & \\ \\

    \theutterance \stepcounter{utterance}  
    & \multicolumn{4}{p{0.6\linewidth}}{
        \cellcolor[rgb]{0.9,0.9,0.9}{
            \makecell[{{p{\linewidth}}}]{
                \texttt{\tiny{[P1$\rangle$GM]}}
                \texttt{PISTA: cuerpo} \\
\texttt{OBJETIVOS: regazo} \\
            }
        }
    }
    & & \\ \\

    \theutterance \stepcounter{utterance}  
    & & & \multicolumn{2}{p{0.3\linewidth}}{
        \cellcolor[rgb]{0.9,0.9,0.9}{
            \makecell[{{p{\linewidth}}}]{
                \texttt{\tiny{[GM$|$GM]}}
                \texttt{cuerpo} \\
            }
        }
    }
    & & \\ \\

    \theutterance \stepcounter{utterance}  
    & & & \multicolumn{2}{p{0.3\linewidth}}{
        \cellcolor[rgb]{0.9,0.9,0.9}{
            \makecell[{{p{\linewidth}}}]{
                \texttt{\tiny{[GM$|$GM]}}
                ['regazo']
            }
        }
    }
    & & \\ \\

    \theutterance \stepcounter{utterance}  
    & & & \multicolumn{2}{p{0.3\linewidth}}{
        \cellcolor[rgb]{0.9,0.9,0.9}{
            \makecell[{{p{\linewidth}}}]{
                \texttt{\tiny{[GM$|$GM]}}
                {'word': 'regazo', 'assignment': 'team'}
            }
        }
    }
    & & \\ \\

    \theutterance \stepcounter{utterance}  
    & & \multicolumn{4}{p{0.6\linewidth}}{
        \cellcolor[rgb]{0.9,0.9,0.9}{
            \makecell[{{p{\linewidth}}}]{
                \texttt{\tiny{[GM$\rangle$P2]}}
                \texttt{La palabra washington fue adivinada correctamente. Tu turno terminó allí.} \\
\texttt{Ahora proporcione otra lista separada por comas de al menos 1 y hasta 1 palabras de la siguiente lista de palabras que mejor se relacionen o estén más estrechamente asociadas con la palabra 'cuerpo'. Recuerde comenzar su respuesta con ADIVINANZA: . Tenga en cuenta: algunas palabras se han eliminado de la lista en comparación con solicitudes anteriores.} \\
\\ 
\texttt{enterrador, aceite, pez, regazo, enfermedad, ciclo, helicóptero, agua, chocolate, prensa, francia, ballena, londres, cara} \\
            }
        }
    }
    & \\ \\

    \theutterance \stepcounter{utterance}  
    & & & \multicolumn{4}{p{0.6\linewidth}}{
        \cellcolor[rgb]{0.9,0.9,0.9}{
            \makecell[{{p{\linewidth}}}]{
                \texttt{\tiny{[GM$\langle$P2]}}
                \texttt{ADIVINANZA: cara} \\
            }
        }
    }
     \\ \\

    \theutterance \stepcounter{utterance}  
    & & & \multicolumn{2}{p{0.3\linewidth}}{
        \cellcolor[rgb]{0.9,0.9,0.9}{
            \makecell[{{p{\linewidth}}}]{
                \texttt{\tiny{[GM$|$GM]}}
                ['cara']
            }
        }
    }
    & & \\ \\

    \theutterance \stepcounter{utterance}  
    & & & \multicolumn{2}{p{0.3\linewidth}}{
        \cellcolor[rgb]{0.9,0.9,0.9}{
            \makecell[{{p{\linewidth}}}]{
                \texttt{\tiny{[GM$|$GM]}}
                {'word': 'cara', 'assignment': 'innocent'}
            }
        }
    }
    & & \\ \\

    \theutterance \stepcounter{utterance}  
    & & & \multicolumn{2}{p{0.3\linewidth}}{
        \cellcolor[rgb]{0.9,0.9,0.9}{
            \makecell[{{p{\linewidth}}}]{
                \texttt{\tiny{[GM$|$GM]}}
                \texttt{cara} \\
            }
        }
    }
    & & \\ \\

    \theutterance \stepcounter{utterance}  
    & & & \multicolumn{2}{p{0.3\linewidth}}{
        \cellcolor[rgb]{0.9,0.9,0.9}{
            \makecell[{{p{\linewidth}}}]{
                \texttt{\tiny{[GM$|$GM]}}
                {'word': 'ciclo', 'assignment': 'opponent'}
            }
        }
    }
    & & \\ \\

    \theutterance \stepcounter{utterance}  
    & & & \multicolumn{2}{p{0.3\linewidth}}{
        \cellcolor[rgb]{0.9,0.9,0.9}{
            \makecell[{{p{\linewidth}}}]{
                \texttt{\tiny{[GM$|$GM]}}
                {'hidden': {'team': ['prensa', 'enterrador', 'pez', 'francia', 'regazo'], 'innocent': ['helicóptero', 'aceite', 'chocolate', 'agua', 'londres'], 'opponent': ['ballena'], 'assassin': ['enfermedad']}, 'revealed': {'team': {'team': ['corneta', 'danza', 'cobre', 'washington'], 'innocent': ['australia', 'cara'], 'opponent': ['olympus'], 'assassin': []}, 'opponent': {'team': [], 'innocent': [], 'opponent': ['conductor', 'ornitorrinco', 'parque', 'pasta', 'mosca', 'ciclo'], 'assassin': []}}}
            }
        }
    }
    & & \\ \\

    \theutterance \stepcounter{utterance}  
    & & \multicolumn{4}{p{0.6\linewidth}}{
        \cellcolor[rgb]{0.9,0.9,0.9}{
            \makecell[{{p{\linewidth}}}]{
                \texttt{\tiny{[P1$\langle$GM]}}
                \texttt{La palabra cara fue adivinada, pero es una palabra inocente. El turno de tu compañero terminó allí.} \\
\texttt{Ahora proporcione otra pista relacionada con algunas de las palabras restantes de su equipo y una lista de las palabras objetivo relacionadas. Recuerde que debe comenzar su pista con 'PISTA: ', poner una nueva línea, y comenzar su lista de palabras objetivo separadas por comas con 'OBJETIVOS: '. Aviso: se han eliminado algunas palabras de las listas con respecto a solicitudes anteriores.} \\
\\ 
\texttt{Las palabras restantes de su equipo son: prensa, enterrador, pez, francia, regazo.} \\
\texttt{Las palabras restantes para su adversario son: ballena, ciclo.} \\
\texttt{Las palabras distractoras restantes son: helicóptero, aceite, chocolate, agua, londres.} \\
\texttt{Las palabras asesinas restantes son: enfermedad.} \\
            }
        }
    }
    & \\ \\

    \theutterance \stepcounter{utterance}  
    & \multicolumn{4}{p{0.6\linewidth}}{
        \cellcolor[rgb]{0.9,0.9,0.9}{
            \makecell[{{p{\linewidth}}}]{
                \texttt{\tiny{[P1$\rangle$GM]}}
                \texttt{PISTA: animal} \\
\texttt{OBJETIVOS: pez} \\
            }
        }
    }
    & & \\ \\

    \theutterance \stepcounter{utterance}  
    & & & \multicolumn{2}{p{0.3\linewidth}}{
        \cellcolor[rgb]{0.9,0.9,0.9}{
            \makecell[{{p{\linewidth}}}]{
                \texttt{\tiny{[GM$|$GM]}}
                \texttt{animal} \\
            }
        }
    }
    & & \\ \\

    \theutterance \stepcounter{utterance}  
    & & & \multicolumn{2}{p{0.3\linewidth}}{
        \cellcolor[rgb]{0.9,0.9,0.9}{
            \makecell[{{p{\linewidth}}}]{
                \texttt{\tiny{[GM$|$GM]}}
                ['pez']
            }
        }
    }
    & & \\ \\

    \theutterance \stepcounter{utterance}  
    & & & \multicolumn{2}{p{0.3\linewidth}}{
        \cellcolor[rgb]{0.9,0.9,0.9}{
            \makecell[{{p{\linewidth}}}]{
                \texttt{\tiny{[GM$|$GM]}}
                {'word': 'pez', 'assignment': 'team'}
            }
        }
    }
    & & \\ \\

    \theutterance \stepcounter{utterance}  
    & & \multicolumn{4}{p{0.6\linewidth}}{
        \cellcolor[rgb]{0.9,0.9,0.9}{
            \makecell[{{p{\linewidth}}}]{
                \texttt{\tiny{[GM$\rangle$P2]}}
                \texttt{La palabra cara fue adivinada, pero es una palabra inocente. Tu turno terminó allí.} \\
\texttt{Ahora proporcione otra lista separada por comas de al menos 1 y hasta 1 palabras de la siguiente lista de palabras que mejor se relacionen o estén más estrechamente asociadas con la palabra 'animal'. Recuerde comenzar su respuesta con ADIVINANZA: . Tenga en cuenta: algunas palabras se han eliminado de la lista en comparación con solicitudes anteriores.} \\
\\ 
\texttt{enterrador, aceite, pez, regazo, enfermedad, helicóptero, agua, chocolate, prensa, francia, ballena, londres} \\
            }
        }
    }
    & \\ \\

    \theutterance \stepcounter{utterance}  
    & & & \multicolumn{4}{p{0.6\linewidth}}{
        \cellcolor[rgb]{0.9,0.9,0.9}{
            \makecell[{{p{\linewidth}}}]{
                \texttt{\tiny{[GM$\langle$P2]}}
                \texttt{ADIVINANZA: ballena} \\
            }
        }
    }
     \\ \\

    \theutterance \stepcounter{utterance}  
    & & & \multicolumn{2}{p{0.3\linewidth}}{
        \cellcolor[rgb]{0.9,0.9,0.9}{
            \makecell[{{p{\linewidth}}}]{
                \texttt{\tiny{[GM$|$GM]}}
                ['ballena']
            }
        }
    }
    & & \\ \\

    \theutterance \stepcounter{utterance}  
    & & & \multicolumn{2}{p{0.3\linewidth}}{
        \cellcolor[rgb]{0.9,0.9,0.9}{
            \makecell[{{p{\linewidth}}}]{
                \texttt{\tiny{[GM$|$GM]}}
                {'word': 'ballena', 'assignment': 'opponent'}
            }
        }
    }
    & & \\ \\

    \theutterance \stepcounter{utterance}  
    & & & \multicolumn{2}{p{0.3\linewidth}}{
        \cellcolor[rgb]{0.9,0.9,0.9}{
            \makecell[{{p{\linewidth}}}]{
                \texttt{\tiny{[GM$|$GM]}}
                \texttt{ballena} \\
            }
        }
    }
    & & \\ \\

    \theutterance \stepcounter{utterance}  
    & & & \multicolumn{2}{p{0.3\linewidth}}{
        \cellcolor[rgb]{0.9,0.9,0.9}{
            \makecell[{{p{\linewidth}}}]{
                \texttt{\tiny{[GM$|$GM]}}
                \texttt{opponent has won} \\
            }
        }
    }
    & & \\ \\

\end{supertabular}
}

\end{document}
